% sections/08_limitations.tex
\section{Limitations}\label{sec:limitations}

The \ISI is subject to the following structural limitations.  These are
inherent to the methodology and cannot be resolved by parameter tuning,
software correction, or improved data quality within the v\,1.0 framework.

\begin{enumerate}[label=\textbf{L\arabic*}.]

  \item \textbf{Re-export and entrepôt masking.}
    The index attributes flows to the immediate bilateral partner, not to
    the country of origin or ultimate destination.  Entrepôt economies
    (notably the Netherlands and Belgium within the EU-27) inflate their
    share in partner-level distributions for goods that merely transit
    through their territory.  This affects axes~3 (semiconductors),
    5~(critical raw materials), and 6~(logistics) most acutely.  No
    harmonised re-export adjustment exists across all source datasets;
    consequently, the \ISI does not strip out re-export flows.

  \item \textbf{Small-economy amplification.}
    Countries with very small absolute import volumes (e.g., CY, MT,
    LU) tend to exhibit high \HHI scores because a single large
    shipment from one partner can dominate the share distribution.
    The resulting concentration score is mechanically correct but may
    overstate structural dependence: a small economy can often switch
    suppliers at lower absolute cost than a large economy with the
    same score.

  \item \textbf{Geographic determinism (Axis~6).}
    Freight flows are partly determined by geography.  Landlocked countries
    (AT, CZ, HU, LU, SK) necessarily concentrate road and rail freight on
    neighbouring partners, producing high logistics \HHI values that
    reflect physical proximity rather than strategic vulnerability.  The
    index does not adjust for geographic proximity or contiguity.

  \item \textbf{No domestic production offset.}
    The \ISI measures concentration of \emph{imports} only.  It does not
    account for domestic production capacity.  A country that imports a
    concentrated bundle of semiconductors but also manufactures a large
    share domestically receives the same axis score as a country with
    identical import patterns but no domestic production.  Import
    concentration is therefore a necessary but not sufficient indicator
    of actual supply vulnerability.

  \item \textbf{Equal weighting arbitrariness.}
    The composite assigns equal weight to all six axes.  This is transparent
    and neutral but may not reflect the relative strategic importance of each
    domain for a specific country.  A small island state's energy-import
    concentration may be existentially significant, while its defence-transfer
    concentration may be economically negligible.  The composite does not
    capture this asymmetry.  Users requiring domain-specific analysis should
    examine axis-level scores directly.

  \item \textbf{Intra-EU trade inclusion.}
    Bilateral flows between EU~member states are included in the share
    vectors.  A country whose imports are concentrated on a single
    EU~partner (e.g., Germany) receives a high \HHI, even though the
    flow occurs within the Single Market and is subject to EU regulatory
    safeguards.  This may overstate effective external dependence for
    countries whose concentration is predominantly intra-EU.  An
    alternative specification---excluding intra-EU flows---would require a
    separate methodology variant and is reserved for future consideration.

  \item \textbf{Single vintage; no time-series comparability.}
    Only \ISI~2024 is materialisable under v\,1.0
    (\cref{sec:data:constraints}).  Cross-temporal comparisons are not
    supported.  Users cannot infer trends, convergence, or divergence from a
    single snapshot.  Time-series extension requires additional source
    vintages and constitutes a future methodology release.

  \item \textbf{No formal verification of the pipeline.}
    The computation pipeline is empirically validated through the test suite
    (\cref{sec:verification}) but has not been subjected to formal
    verification (proof-carrying code, model checking, or certified
    compilation).  The determinism invariants (\cref{sec:det:invariants}) are
    contractual obligations enforced by testing, not by mathematical proof.
    Residual implementation defects that escape the test suite remain
    possible.

\end{enumerate}

These limitations are permanent features of v\,1.0.  They are documented
here to enable informed use of the index and to prevent over-interpretation
of the published scores.